\section{Results}

The results are presented according to the diagnostic questions addressed by the analysis. Figure~\ref{fig:workflow} summarizes the workflow: public molecular property datasets were featurized consistently, evaluated under three split protocols, and examined using target-shift, scaffold, chemical-similarity, performance-gap, and ranking-stability diagnostics.

\begin{figure}[H]
\centering
\includegraphics[width=\textwidth]{figure1_split_diagnostic_workflow.png}
\caption{Overview of the split-diagnostic audit. (A) Molecular property datasets are standardized, featurized, split, modelled, and interpreted through a diagnostic workflow. (B) The benchmark spans classification and regression datasets of different sizes. (C) Each split is audited through target-shift, scaffold, chemical-similarity, and model-stability diagnostics. (D) Random, ordinary scaffold, and target-balanced scaffold splits emphasize different benchmark questions rather than forming a single universal ranking of difficulty.}
\label{fig:workflow}
\end{figure}

\subsection{Scaffold separation alone does not guarantee a clean benchmark split}

Ordinary scaffold splitting removed direct train/test scaffold overlap across the tested datasets, but it also changed other properties of the benchmark. In particular, scaffold separation can be accompanied by target distribution shift and scaffold concentration. ESOL and FreeSolv illustrate this issue most clearly: ordinary scaffold splitting holds out very few scaffold groups and produces large target mean gaps. Target-balanced scaffold splitting reduces this problem in several datasets, but FreeSolv remains constrained by small dataset size and a dominant scaffold group.

Figure~\ref{fig:target-shift} shows the target-shift component directly. Random splits have small target gaps because molecules are distributed without scaffold-level constraints and stratification is used where applicable. Ordinary scaffold splits can be more shifted because complete scaffold groups are held out. Target-balanced scaffold splitting often reduces the gap, although strong scaffold concentration can still limit the achievable balance.

\begin{figure}[H]
\centering
\includegraphics[width=0.94\textwidth]{figure3_target_shift.png}
\caption{Target distribution shift under random, ordinary scaffold, and target-balanced scaffold splitting. Target-balanced scaffold splitting reduces target mean gaps in several datasets, but small or scaffold-concentrated datasets can remain constrained.}
\label{fig:target-shift}
\end{figure}

\subsection{Performance gaps shrink in some datasets after target balancing, but not universally}

Figure~\ref{fig:gap} summarizes the primary generalization-gap result. In several datasets, especially ESOL for tree-based models, the gap between random and ordinary scaffold evaluation decreases after target balancing. This pattern indicates that part of the ordinary scaffold gap can be associated with target shift rather than scaffold separation alone. The effect is not uniform across tasks or models. In HIV, balanced scaffold gaps are larger than ordinary scaffold gaps across models, and in FreeSolv different models respond differently to balancing. These results indicate that target balancing changes the evaluation condition rather than producing a universally easier or harder split.

\begin{figure}[H]
\centering
\includegraphics[width=\textwidth]{figure2_generalization_gap.png}
\caption{Generalization gap under ordinary scaffold and target-balanced scaffold splitting. Classification gaps are random ROC-AUC minus scaffold ROC-AUC. Regression gaps are scaffold RMSE minus random RMSE. Positive values indicate worse performance under the scaffold-based split than under the random split.}
\label{fig:gap}
\end{figure}

\subsection{Random splits are chemically closer to the training set}

Figure~\ref{fig:tanimoto} shows that random split test molecules generally have higher maximum Tanimoto similarity to training molecules than scaffold-based test molecules. This supports the interpretation that random-split performance can reflect interpolation within chemically similar regions. Ordinary scaffold splitting usually produces lower test-to-train similarity, while target-balanced scaffold splitting often lies between random and ordinary scaffold splitting. This intermediate behavior separates, at least partially, target-shift control from maximal chemical dissimilarity.

\begin{figure}[H]
\centering
\includegraphics[width=0.94\textwidth]{figure4_tanimoto_similarity.png}
\caption{Mean maximum test-to-train Tanimoto similarity by split strategy. Random splits generally place test molecules closer to the training set than scaffold-based splits, making aggregate random-split scores difficult to interpret as scaffold-level generalization.}
\label{fig:tanimoto}
\end{figure}

\subsection{Split diagnostics change model comparison}

The appendix ranking analysis shows that split choice can change not only absolute performance estimates but also which model appears best. This matters because benchmark studies often use model rankings to support claims about method superiority. If the best-performing model changes when the split changes, then a leaderboard-style conclusion is incomplete without the corresponding split diagnostics.

Taken together, the results show that the split protocol changes several aspects of benchmark interpretation. Random splitting often produces chemically closer test sets. Ordinary scaffold splitting removes scaffold overlap but can introduce target shift and scaffold concentration. Target-balanced scaffold splitting helps separate target-shift effects from scaffold-generalization effects, but does not eliminate all dataset- and model-dependent behavior. Full numerical values are retained in the generated result tables and appendix materials to preserve reproducibility.
